% [VERIFY] related work traces to: related_work.md RW01-RW12 + briefer citations; typed edges preserved
\section{Related Work}
\label{sec:related}

\noindent\textbf{Residual representations.}
The idea of representing a signal by its residual with respect to a reference has a long
history in vision and numerical computation. In image retrieval and classification,
encoding residual vectors---for example, by quantizing residuals relative to a
dictionary---has proven more effective than encoding original vectors, as in product
quantization and related Fisher/VLAD encodings~\citep{jegou2011pq}. In low-level vision
and graphics, multigrid and hierarchical-basis preconditioning solve partial
differential equations far faster by reformulating the system in terms of residual
variables at multiple scales~\citep{briggs2000multigrid,szeliski2006preconditioning}.
These methods establish a recurring principle---that a good reformulation or
preconditioning can simplify optimization---but operate on hand-designed
representations. We carry the principle into deep learning by making every few stacked
layers fit a residual function, so that the reformulation is learned rather than
engineered.

\noindent\textbf{Shortcut connections and gating.}
Skip and shortcut connections predate deep residual learning, appearing in early
multilayer-perceptron design and in the auxiliary-classifier branches used to combat
vanishing gradients~\citep{lee2014dsn}. The closest prior method is the highway
network~\citep{srivastava2015highway,srivastava2015trainingverydeep}, whose shortcuts
are governed by data-dependent gates with their own parameters. Gating couples the
shortcut to extra parameters and, more importantly, can ``close'': when a gate
approaches zero the layer reverts to a non-residual transformation, and highway networks
have not shown accuracy gains beyond roughly one hundred layers. We adopt the general
notion of a skip connection but reject the gating mechanism, using instead a
parameter-free identity shortcut that is never closed and never adds parameters, which
is what allows depth to be increased without bound on the optimization side.

\noindent\textbf{Deep architectures and training of plain networks.}
A second line of work scales plain feed-forward depth directly. VGG nets establish the
design philosophy of stacking small $3\times3$ convolutions under simple width
rules~\citep{simonyan2015vgg}, and we adopt these rules but reach far greater depth at
much lower complexity. Inception/GoogLeNet networks add depth through a branching
topology with internal shortcut-and-merge structure~\citep{szegedy2015googlenet}, which
we use only as a comparison point rather than as an architectural ingredient. The
ingredients that make very deep plain training feasible at all---batch
normalization~\citep{ioffe2015batchnorm}, variance-preserving initialization for
rectifiers~\citep{he2015prelu,glorot2010xavier}, and rectified linear
units~\citep{nair2010relu}---are imported wholesale and applied identically to plain and
residual variants, so that any difference between them is attributable to the residual
reformulation and not to the training recipe. These techniques let tens-of-layer plain
nets converge, but, as the long-recognized difficulty of training deep
networks~\citep{bengio1994longterm} would predict, they do not by themselves remove the
\degprob, which is precisely the gap residual learning closes.

\noindent\textbf{Object detection frameworks.}
For the transfer experiments we build on the Faster R-CNN detection
framework~\citep{ren2015fasterrcnn} together with the Fast R-CNN region
classifier~\citep{girshick2015fastrcnn} and its R-CNN predecessor~\citep{girshick2014rcnn},
keeping the detection pipeline---region-proposal network, RoI pooling, and four-step
alternating training---fixed. To insert a fully-convolutional residual backbone that
lacks hidden fully-connected layers, we follow the strategy of sharing full-image
convolutional features and treating a late convolutional stage as the per-region
classifier~\citep{ren2015noc}. Holding the framework constant and changing only the
backbone isolates the contribution of the learned representation, and we evaluate it on
the standard PASCAL VOC~\citep{everingham2010pascalvoc}, ImageNet~\citep{russakovsky2015imagenet},
and COCO~\citep{lin2014coco} benchmarks, with the CIFAR-10 analysis using the standard
tiny-image dataset~\citep{krizhevsky2009cifar} and augmentation conventions established
in prior recognition work~\citep{krizhevsky2012alexnet}. Baselines from this literature,
including network-in-network~\citep{lin2013nin}, maxout~\citep{goodfellow2013maxout}, and
thin deep nets~\citep{romero2015fitnets}, anchor our CIFAR-10 comparison, and fully
convolutional inference~\citep{long2015fcn} underlies the multi-scale testing protocol.
